\documentclass[aspectratio=169]{beamer}
\usetheme{UFS}
\usepackage[utf8]{inputenc}
\usepackage[portuguese]{babel}
\usepackage{amsmath,amssymb}
\usepackage{graphicx}
\usepackage{booktabs}
\usepackage{listings}
\usepackage{xcolor}
\usepackage{tikz}

% ── VS Code Light+ palette ──
\definecolor{bgcolor}{HTML}{FFFFFF}
\definecolor{linecolor}{HTML}{DDDDDD}
\definecolor{numcolor}{HTML}{999999}
\definecolor{kwcolor}{HTML}{0000FF}
\definecolor{strcolor}{HTML}{A31515}
\definecolor{cmtcolor}{HTML}{008000}
\definecolor{fncolor}{HTML}{795E26}
\definecolor{typecolor}{HTML}{267F99}

\lstdefinestyle{modernstyle}{
  language=Python,
  backgroundcolor=\color{bgcolor},
  rulecolor=\color{linecolor},
  numberstyle=\tiny\color{numcolor},
  keywordstyle=\color{kwcolor}\bfseries,
  stringstyle=\color{strcolor},
  commentstyle=\color{cmtcolor}\itshape,
  emphstyle=\color{fncolor},
  emphstyle={[2]\color{typecolor}},
  basicstyle=\ttfamily\scriptsize,
  breaklines=true,
  frame=single,
  numbers=left,
  numbersep=8pt,
  tabsize=4,
  showstringspaces=false,
  columns=flexible,
  literate={á}{{\'a}}1 {ã}{{\~a}}1 {é}{{\'e}}1 {ê}{{\^e}}1 {í}{{\'i}}1
           {ó}{{\'o}}1 {ô}{{\^o}}1 {ú}{{\'u}}1 {ç}{{\c{c}}}1
           {Á}{{\'A}}1 {É}{{\'E}}1 {Í}{{\'I}}1 {Ó}{{\'O}}1 {Ú}{{\'U}}1,
  emph={receber_pagamento,ver_saldo,preco,total,__init__,__str__,__repr__,
        __eq__,__lt__,__hash__,__len__,__add__,__sub__,__mul__,simplificar,
        setter,deleter},
  emph={[2]int,float,str,bool,list,dict,set,tuple,property,classmethod,
        staticmethod,ValueError,AttributeError,TypeError},
}
\lstset{style=modernstyle}

\title{COMP0496 -- Programação A \\ Programação Orientada a Objetos --- Parte 2}
\subtitle{Encapsulamento e Métodos Especiais}
\author{Prof.\ Dr.\ André Yoshiaki Kashiwabara}
\institute{DCOMP -- Universidade Federal de Sergipe\\\texttt{andreyoshiaki@dcomp.ufs.br}}
\date{}

\begin{document}

\begin{frame}
\titlepage
\end{frame}

\begin{frame}{Roteiro}
\tableofcontents
\end{frame}

% ================================================================
\section{Revisão --- O Problema do Estado Inválido}
% ================================================================

\begin{frame}[fragile]{Relembrando: a classe \texttt{Lanche}}
Na Aula~1, criamos classes para o food truck \textbf{O Podrão}.

\begin{lstlisting}
class Lanche:
    def __init__(self, nome, preco):
        self.nome = nome
        self.preco = preco

xtudo = Lanche("X-Tudo", 18.50)
\end{lstlisting}

Tudo funciona bem --- até alguém fazer algo inesperado.
\end{frame}

\begin{frame}[fragile]{O problema: estado inválido}
\begin{lstlisting}
xtudo = Lanche("X-Tudo", 18.50)
xtudo.preco = -999      # nenhum erro!
xtudo.preco = "gratis"  # nenhum erro!
print(xtudo.preco)      # "gratis"
\end{lstlisting}

\begin{alertblock}{Problema}
Qualquer código externo pode colocar o objeto em um estado que não faz sentido.
Um preço negativo ou textual viola as regras do nosso domínio.
\end{alertblock}

\medskip
Como proteger os dados internos de um objeto?
\end{frame}

% ================================================================
\section{Encapsulamento}
% ================================================================

\begin{frame}{Analogia: a caixa registradora}
\begin{columns}[T]
\column{0.55\textwidth}
Pense na caixa registradora do Podrão:
\begin{itemize}
    \item O \textbf{visor} mostra o total (acesso público)
    \item A \textbf{gaveta} só abre com a chave do operador (acesso restrito)
    \item Ninguém coloca a mão diretamente no dinheiro
\end{itemize}

\column{0.42\textwidth}
\begin{block}{Encapsulamento}
Esconder os detalhes internos de um objeto e expor apenas uma
interface controlada.
\end{block}
\end{columns}
\end{frame}

\begin{frame}{Convenções Python: público, protegido, privado}
Python \textbf{não tem} modificadores de acesso como Java.
Usa-se \textbf{convenções de nomenclatura}:

\medskip
\begin{center}
\begin{tabular}{lll}
\toprule
\textbf{Nome} & \textbf{Convenção} & \textbf{Significado} \\
\midrule
\texttt{nome}    & público   & Acesso livre \\
\texttt{\_nome}  & protegido & ``Use com cuidado'' \\
\texttt{\_\_nome}& privado   & Name mangling (dificulta acesso) \\
\bottomrule
\end{tabular}
\end{center}

\medskip
\begin{exampleblock}{Filosofia Python}
``We're all consenting adults here'' --- a linguagem confia no programador,
mas oferece mecanismos para sinalizar intenção.
\end{exampleblock}
\end{frame}

\begin{frame}[fragile]{Name mangling: como funciona}
Atributos com \texttt{\_\_} (duplo sublinhado) sofrem \textbf{name mangling}:

\begin{lstlisting}
class CaixaRegistradora:
    def __init__(self):
        self.__saldo = 0.0

    def receber_pagamento(self, valor):
        if valor > 0:
            self.__saldo += valor

    def ver_saldo(self):
        return self.__saldo
\end{lstlisting}
\end{frame}

\begin{frame}[fragile]{Name mangling: acesso externo}
\begin{lstlisting}
caixa = CaixaRegistradora()
caixa.receber_pagamento(25.00)
print(caixa.ver_saldo())   # 25.0

print(caixa.__saldo)
# AttributeError: 'CaixaRegistradora' object
#   has no attribute '__saldo'
\end{lstlisting}

\begin{alertblock}{Name mangling por dentro}
O Python renomeia \texttt{\_\_saldo} para \texttt{\_CaixaRegistradora\_\_saldo}.
\end{alertblock}

\begin{lstlisting}
# Funciona, mas NUNCA faca isso!
print(caixa._CaixaRegistradora__saldo)  # 25.0
\end{lstlisting}
\end{frame}

% ================================================================
\section{@property --- Acesso Controlado}
% ================================================================

\begin{frame}[fragile]{O problema dos getters e setters}
Em Java, usamos \texttt{getPreco()} e \texttt{setPreco()}.
Em Python, isso fica verboso e pouco natural:

\begin{lstlisting}
# Estilo Java em Python (evite!)
lanche.get_preco()
lanche.set_preco(20.0)

# Estilo Pythonico (queremos isto)
lanche.preco        # leitura
lanche.preco = 20.0 # escrita com validacao
\end{lstlisting}

A solução: \texttt{@property}.
\end{frame}

\begin{frame}[fragile]{\texttt{@property}: métodos disfarçados de atributos}
\begin{lstlisting}
class Lanche:
    def __init__(self, nome, preco):
        self._nome = nome
        self.preco = preco  # usa o setter!

    @property
    def preco(self):
        return self._preco

    @preco.setter
    def preco(self, valor):
        if valor < 0:
            raise ValueError("Preco negativo!")
        if valor > 500:
            raise ValueError("Preco acima do limite!")
        self._preco = valor
\end{lstlisting}
\end{frame}

\begin{frame}[fragile]{Usando a property}
\begin{lstlisting}
xtudo = Lanche("X-Tudo", 18.50)
print(xtudo.preco)   # 18.5 (chama o getter)

xtudo.preco = 22.0   # OK (chama o setter)
xtudo.preco = -10     # ValueError: Preco negativo!
xtudo.preco = 1000    # ValueError: Preco acima do limite!
\end{lstlisting}

\begin{block}{Vantagem}
O código externo usa sintaxe de atributo, mas internamente
há validação. Mudança transparente.
\end{block}
\end{frame}

\begin{frame}[fragile]{Property somente-leitura}
Use \texttt{@property} \textbf{sem} definir o setter:

\begin{lstlisting}
class Pedido:
    def __init__(self):
        self._itens = []

    def adicionar(self, lanche, qtd=1):
        self._itens.append((lanche, qtd))

    @property
    def total(self):
        return sum(l.preco * q for l, q in self._itens)
\end{lstlisting}

\begin{lstlisting}
pedido = Pedido()
pedido.adicionar(xtudo, 2)
print(pedido.total)     # 37.0
pedido.total = 0        # AttributeError!
\end{lstlisting}
\end{frame}

\begin{frame}[fragile]{Armadilha: recursão infinita}
\begin{alertblock}{Cuidado!}
Não use o mesmo nome para a property e o atributo interno.
\end{alertblock}

\begin{lstlisting}
class Errado:
    @property
    def preco(self):
        return self.preco  # chama a si mesmo!

    @preco.setter
    def preco(self, valor):
        self.preco = valor  # chama a si mesmo!
    # RecursionError!
\end{lstlisting}

\textbf{Solução:} armazene em \texttt{self.\_preco} (com prefixo).
\end{frame}

% ================================================================
\section{Métodos Especiais (Dunders)}
% ================================================================

\begin{frame}{O que são dunders?}
Métodos com duplo sublinhado: \texttt{\_\_nome\_\_}

São chamados automaticamente pelo Python em situações específicas:

\medskip
\begin{center}
\begin{tabular}{ll}
\toprule
\textbf{Você escreve} & \textbf{Python chama} \\
\midrule
\texttt{print(obj)}      & \texttt{obj.\_\_str\_\_()} \\
\texttt{repr(obj)}       & \texttt{obj.\_\_repr\_\_()} \\
\texttt{obj1 == obj2}    & \texttt{obj1.\_\_eq\_\_(obj2)} \\
\texttt{obj1 < obj2}     & \texttt{obj1.\_\_lt\_\_(obj2)} \\
\texttt{len(obj)}        & \texttt{obj.\_\_len\_\_()} \\
\texttt{hash(obj)}       & \texttt{obj.\_\_hash\_\_()} \\
\texttt{obj1 + obj2}     & \texttt{obj1.\_\_add\_\_(obj2)} \\
\bottomrule
\end{tabular}
\end{center}
\end{frame}

\begin{frame}[fragile]{\texttt{\_\_str\_\_} vs \texttt{\_\_repr\_\_}}
\begin{lstlisting}
class Lanche:
    def __init__(self, nome, preco):
        self._nome = nome
        self._preco = preco

xtudo = Lanche("X-Tudo", 18.5)
print(xtudo)
# <__main__.Lanche object at 0x7f3b2c>  (feio!)
\end{lstlisting}

Agora com dunders:
\begin{lstlisting}
    def __str__(self):
        return f"{self._nome} - R${self._preco:.2f}"

    def __repr__(self):
        return f"Lanche('{self._nome}', {self._preco})"
\end{lstlisting}

\begin{lstlisting}
print(xtudo)        # X-Tudo - R$18.50  (amigavel)
repr(xtudo)         # Lanche('X-Tudo', 18.5) (tecnico)
\end{lstlisting}
\end{frame}

\begin{frame}[fragile]{\texttt{\_\_eq\_\_}: igualdade por valor}
Por padrão, dois objetos com os mesmos dados são \textbf{diferentes}:

\begin{lstlisting}
a = Lanche("X-Tudo", 18.5)
b = Lanche("X-Tudo", 18.5)
print(a == b)  # False (compara identidade!)
\end{lstlisting}

Solução:
\begin{lstlisting}
    def __eq__(self, other):
        if not isinstance(other, Lanche):
            return NotImplemented
        return (self._nome == other._nome
                and self._preco == other._preco)
\end{lstlisting}

\begin{lstlisting}
print(a == b)  # True
\end{lstlisting}
\end{frame}

\begin{frame}[fragile]{\texttt{\_\_lt\_\_}: ordenação com \texttt{sorted()}}
\begin{lstlisting}
    def __lt__(self, other):
        if not isinstance(other, Lanche):
            return NotImplemented
        return self._preco < other._preco
\end{lstlisting}

\begin{lstlisting}
cardapio = [
    Lanche("X-Tudo", 18.5),
    Lanche("Misto",   8.0),
    Lanche("Burgao", 22.0),
]
for l in sorted(cardapio):
    print(l)
# Misto - R$8.00
# X-Tudo - R$18.50
# Burgao - R$22.00
\end{lstlisting}
\end{frame}

\begin{frame}[fragile]{\texttt{\_\_hash\_\_}: objetos em conjuntos e dicionários}
Ao definir \texttt{\_\_eq\_\_}, o Python desabilita \texttt{\_\_hash\_\_}:

\begin{lstlisting}
lanches = {Lanche("X-Tudo", 18.5)}
# TypeError: unhashable type: 'Lanche'
\end{lstlisting}

Solução:
\begin{lstlisting}
    def __hash__(self):
        return hash((self._nome, self._preco))
\end{lstlisting}

\begin{alertblock}{Regra de ouro}
Se \texttt{a == b}, então \texttt{hash(a) == hash(b)}.
Use os mesmos atributos em ambos os métodos.
\end{alertblock}
\end{frame}

\begin{frame}[fragile]{\texttt{\_\_len\_\_}: tamanho de coleções}
\begin{lstlisting}
class Pedido:
    def __init__(self):
        self._itens = []

    def adicionar(self, lanche, qtd=1):
        self._itens.append((lanche, qtd))

    def __len__(self):
        return sum(q for _, q in self._itens)
\end{lstlisting}

\begin{lstlisting}
pedido = Pedido()
pedido.adicionar(Lanche("X-Tudo", 18.5), 2)
pedido.adicionar(Lanche("Misto", 8.0), 1)
print(len(pedido))  # 3
\end{lstlisting}
\end{frame}

\begin{frame}{Tabela resumo: principais dunders}
\begin{center}
\scriptsize
\begin{tabular}{lll}
\toprule
\textbf{Método} & \textbf{Operação} & \textbf{Quando usar} \\
\midrule
\texttt{\_\_init\_\_}  & Construtor          & Sempre \\
\texttt{\_\_str\_\_}   & \texttt{print()}, \texttt{str()} & Saída amigável \\
\texttt{\_\_repr\_\_}  & \texttt{repr()}, terminal & Debug, logs \\
\texttt{\_\_eq\_\_}    & \texttt{==}         & Comparação por valor \\
\texttt{\_\_lt\_\_}    & \texttt{<}, \texttt{sorted()} & Ordenação \\
\texttt{\_\_hash\_\_}  & \texttt{set}, \texttt{dict} & Após definir \texttt{\_\_eq\_\_} \\
\texttt{\_\_len\_\_}   & \texttt{len()}      & Coleções \\
\texttt{\_\_add\_\_}   & \texttt{+}          & Soma/concatenação \\
\texttt{\_\_sub\_\_}   & \texttt{-}          & Subtração \\
\texttt{\_\_mul\_\_}   & \texttt{*}          & Multiplicação \\
\bottomrule
\end{tabular}
\end{center}
\end{frame}

% ================================================================
\section{Prática --- Classe Fracao}
% ================================================================

\begin{frame}{Classe \texttt{Fracao}: especificação}
Vamos construir uma classe completa que usa \textbf{tudo} que aprendemos:

\begin{itemize}
    \item Validação no construtor (denominador $\neq 0$)
    \item Simplificação automática via MDC
    \item Properties somente-leitura
    \item \texttt{\_\_str\_\_}, \texttt{\_\_repr\_\_}
    \item \texttt{\_\_eq\_\_}, \texttt{\_\_lt\_\_}, \texttt{\_\_hash\_\_}
    \item \texttt{\_\_add\_\_}, \texttt{\_\_sub\_\_}, \texttt{\_\_mul\_\_}
\end{itemize}
\end{frame}

\begin{frame}[fragile]{Classe \texttt{Fracao}: construtor e simplificação}
\begin{lstlisting}
from math import gcd

class Fracao:
    def __init__(self, num, den=1):
        if den == 0:
            raise ValueError("Denominador nao pode ser zero!")
        # sinal sempre no numerador
        if den < 0:
            num, den = -num, -den
        divisor = gcd(abs(num), den)
        self._num = num // divisor
        self._den = den // divisor

    @property
    def num(self):
        return self._num

    @property
    def den(self):
        return self._den
\end{lstlisting}
\end{frame}

\begin{frame}[fragile]{Classe \texttt{Fracao}: representação}
\begin{lstlisting}
    def __str__(self):
        if self._den == 1:
            return str(self._num)
        return f"{self._num}/{self._den}"

    def __repr__(self):
        return f"Fracao({self._num}, {self._den})"
\end{lstlisting}

\begin{lstlisting}
print(Fracao(4, 6))   # 2/3
print(Fracao(6, 3))   # 2
repr(Fracao(4, 6))    # Fracao(2, 3)
\end{lstlisting}
\end{frame}

\begin{frame}[fragile]{Classe \texttt{Fracao}: comparação e hash}
\begin{lstlisting}
    def __eq__(self, other):
        if not isinstance(other, Fracao):
            return NotImplemented
        return (self._num == other._num
                and self._den == other._den)

    def __lt__(self, other):
        if not isinstance(other, Fracao):
            return NotImplemented
        return self._num * other._den < other._num * self._den

    def __hash__(self):
        return hash((self._num, self._den))
\end{lstlisting}
\end{frame}

\begin{frame}[fragile]{Classe \texttt{Fracao}: operações aritméticas}
\begin{lstlisting}
    def __add__(self, other):
        if not isinstance(other, Fracao):
            return NotImplemented
        return Fracao(self._num * other._den
                      + other._num * self._den,
                      self._den * other._den)

    def __sub__(self, other):
        if not isinstance(other, Fracao):
            return NotImplemented
        return Fracao(self._num * other._den
                      - other._num * self._den,
                      self._den * other._den)

    def __mul__(self, other):
        if not isinstance(other, Fracao):
            return NotImplemented
        return Fracao(self._num * other._num,
                      self._den * other._den)
\end{lstlisting}
\end{frame}

\begin{frame}[fragile]{Classe \texttt{Fracao}: demonstração}
\begin{lstlisting}
a = Fracao(1, 2)
b = Fracao(1, 3)

print(a + b)          # 5/6
print(a - b)          # 1/6
print(a * b)          # 1/6
print(a == Fracao(2, 4))  # True (simplificado!)
\end{lstlisting}

\begin{lstlisting}
fracs = [Fracao(3,4), Fracao(1,2), Fracao(2,3)]
print(sorted(fracs))  # [1/2, 2/3, 3/4]
\end{lstlisting}

\begin{lstlisting}
# set() funciona porque temos __eq__ e __hash__
s = {Fracao(1,2), Fracao(2,4), Fracao(3,6)}
print(s)              # {1/2}
\end{lstlisting}
\end{frame}

% ================================================================
\section{Resumo}
% ================================================================

\begin{frame}{Resumo da Aula}
\begin{center}
\begin{tabular}{ll}
\toprule
\textbf{Conceito} & \textbf{Ferramenta Python} \\
\midrule
Esconder dados        & \texttt{\_\_atributo} (name mangling) \\
Acesso controlado     & \texttt{@property} / \texttt{@x.setter} \\
Representação textual & \texttt{\_\_str\_\_} / \texttt{\_\_repr\_\_} \\
Igualdade por valor   & \texttt{\_\_eq\_\_} + \texttt{\_\_hash\_\_} \\
Ordenação             & \texttt{\_\_lt\_\_} \\
Operadores            & \texttt{\_\_add\_\_}, \texttt{\_\_sub\_\_}, \texttt{\_\_mul\_\_} \\
\bottomrule
\end{tabular}
\end{center}
\end{frame}

\begin{frame}{Armadilhas comuns}
\begin{alertblock}{Cuidado!}
\begin{enumerate}
    \item Property com mesmo nome do atributo interno $\Rightarrow$ recursão infinita
    \item Definir \texttt{\_\_eq\_\_} sem \texttt{\_\_hash\_\_} $\Rightarrow$ objetos não entram em \texttt{set}/\texttt{dict}
    \item Acessar \texttt{\_CaixaRegistradora\_\_saldo} diretamente $\Rightarrow$ quebra encapsulamento
    \item Esquecer \texttt{isinstance} check nos dunders $\Rightarrow$ erros com tipos mistos
\end{enumerate}
\end{alertblock}
\end{frame}

\begin{frame}{Próxima aula: Herança e Polimorfismo}
\begin{block}{Aula 3 --- Preview}
\begin{itemize}
    \item Herança: \texttt{class XTudo(Lanche)}
    \item Polimorfismo: mesmo método, comportamentos diferentes
    \item Composição vs.\ herança: quando usar cada um
    \item O food truck cresce: cardápio com subclasses
\end{itemize}
\end{block}

\bigskip
\centering
\Large Dúvidas?
\end{frame}

\end{document}
